% !TEX root = ./main.tex

\chapter{系统需求分析与总体设计}

多模态AIGC安全分析平台需要将图像、视频、生成模型、检测模型和外部审核接口组织为统一系统。仅有单个模型或单个页面并不能支撑完整的安全分析任务，平台还需要明确输入形式、任务边界、结果结构、模块职责和部署约束。本章围绕系统建设目标、系统需求、总体架构和业务流程展开分析，将课题目标以及工程约束共同转化为可实现的系统设计方案。

\section{系统建设目标}

本文所研究的系统并不是一个单纯的模型演示页面，而是一个围绕AIGC安全分析任务构建的多模态内容检测与分析平台。该平台的建设目标可以概括为“面向图像与视频内容，形成可接入、可分析、可展示、可扩展的AIGC安全研究原型”。围绕这一目标，系统设计需要同时满足研究验证、工程演示和后续扩展三个层面的要求。

从研究验证角度看，平台应能够承载虚假内容检测与不安全内容审核两类核心任务，并支持本地模型与外部接口的统一接入。这样做的原因在于，本文的目标并非局限于调用某一个现成接口，而是希望在统一平台中比较不同检测能力的适用场景、组织方式和可扩展性。因此，系统必须具备多源模型管理与统一结果表示能力。

从工程演示角度看，平台需要具备完整的交互闭环，包括内容输入、任务触发、结果返回、风险展示和历史追踪等环节。只有在业务链路上形成闭环，系统才能既服务于实验验证，又能够支撑功能展示和持续扩展。换言之，平台不仅要“能跑通模型”，还要“能被用户使用”，这也是本文区别于纯算法复现工作的一个重要特征。

从持续扩展角度看，系统应具备灵活的模块替换能力。AIGC安全研究处于快速发展阶段，无论是生成模型还是检测方法都在持续迭代。如果平台在设计之初就将不同模型、不同输入模态和不同结果结构紧密耦合，模型替换和能力扩展成本会非常高。因此，系统需要在架构层面保持较强的松耦合特征，使前端界面、后端代理、模型服务和外部接口之间能够通过统一协议协同工作。

基于上述分析，本文将系统设计目标进一步细化为以下几点：第一，支持图像与视频两类输入媒介，并兼容文件上传和链接输入两种常见方式；第二，支持虚假内容检测和不安全内容审核两类安全任务；第三，支持本地模型与外部接口协同工作，形成多源能力统一接入框架；第四，支持检测结果可视化、风险分级和历史管理；第五，支持内容生成与检测联动，为样本构建与实验验证提供平台基础。围绕这些目标展开的需求分析，将直接决定系统总体架构与模块划分方式。

\section{系统需求分析}

系统需求需要直接服务于后续架构设计。本文不再将需求拆成大量孤立条目，而是围绕平台必须解决的四个核心问题展开：如何接收图像与视频，如何完成真假检测与安全审核，如何形成生成样本的送检闭环，以及如何保证系统具备继续扩展和稳定部署的工程基础。

\subsection{图像与视频输入需求}

平台输入对象限定为图像和视频两类媒介，输入来源包括本地上传和URL输入两种形式。该部分的核心不是增加复杂入口，而是把媒介类型、来源方式和检测目标在输入阶段确定下来，使后端能够根据统一参数完成任务路由。

对图像输入，系统需要支持文件预览、格式校验以及Base64或URL形式的后端传递；对视频输入，重点在于文件大小限制、播放预览、远程地址透传和长耗时任务状态提示。由于视频处理成本明显高于图像，系统不应把二者完全放入同一套同步处理逻辑，而应在输入阶段完成媒介区分。

因此，输入模块至少应输出三个稳定字段：媒介类型、输入来源和任务类型。前端负责避免用户进入错误流程，后端负责参数校验和异常拦截。这样既能减少后续检测逻辑中的分支混乱，也为接口安全和错误定位提供基础。

\subsection{内容检测与安全审核}

平台的检测能力分为两条主线：虚假内容检测关注内容是否由AI生成或经过伪造处理，安全审核关注内容是否存在违规风险。二者不能合并为同一个判断结果，但可以共用统一的输入结构、结果字段和风险展示组件。

虚假内容检测需要覆盖图像和视频。图像侧以本地模型复现为基础，同时保留外部多模态接口作为开放场景补充；视频侧由于本地模型部署和推理成本较高，主要依赖外部接口完成真假判别。检测结果应包括真假结论、置信度、模型来源和必要原因说明，避免只给出单一标签而缺乏解释和记录价值。

安全审核侧重点不同，核心输出是安全结论、风险标签、处置建议和原因说明。图像审核可直接分析静态内容，视频审核则需要支持更长耗时和更复杂的接口状态。因此，安全审核模块既要和虚假检测共享结果展示方式，也要保留对视频异步任务的处理能力。

从系统角度看，本节最关键的约束是结果归一化。无论底层来自本地模型还是外部平台，前端都应接收统一字段，并根据任务耗时区分同步返回和异步轮询。这样可以避免每接入一个新模型就重写页面逻辑。

\subsection{内容生成、送检与结果管理}

内容生成模块在本文中主要承担样本构建作用，而不是单独作为生成模型展示平台。系统需要支持文生图、文生视频、图生视频以及若干Deepfake相关任务，使平台能够产生可控样本，并立即进入检测或审核流程。

生成结果需要保存为平台样本，并记录媒体地址、来源模块、生成模型、提示词、创建时间和检测状态等必要元数据。用户从生成结果发起“AI生成检测”或“敏感内容检测”时，系统应自动完成样本入库、页面跳转和检测输入填充，减少下载后再上传的重复操作。

检测记录也需要持久化保存。记录字段应覆盖检测类型、媒体类型、检测结论、置信度或风险分数、模型来源、风险标签、原始结果和关联样本编号。当样本被删除时，检测记录仍应保留必要快照，以便后续复核。这样生成、送检和记录管理才能形成闭环。

\subsection{系统扩展与工程质量}

平台后续可能继续接入新的生成模型、检测模型和审核接口，因此扩展能力是需求阶段必须明确的约束。系统应把模型差异限制在后端适配层，前端尽量依赖稳定的数据结构和展示组件。

工程质量主要体现在四个方面：第一，前端、后端和模型服务职责清晰；第二，检测结果、生成样本和历史记录使用稳定结构；第三，图像类短任务支持同步返回，视频和生成类长任务支持状态提示或轮询；第四，接口异常、模型超时和参数错误能够被定位和提示。

安全与部署方面，外部接口密钥必须保留在后端，上传内容需要进行基本校验，本地模型服务与后端代理之间的调用关系需要明确。系统还应具备本地开发和服务器演示两种运行方式，在依赖管理、服务启动和端口配置上保持可迁移性。

\section{系统架构设计}

\subsection{总体架构设计思路}

基于前述需求分析，本文将系统总体架构设计为前端展示层、后端代理层、模型推理与外部能力层三部分。之所以采用三层分离架构，而非将前端直接连接模型或将所有逻辑堆叠在单一服务中，主要基于以下考虑：其一，前端需要保持对多源能力的无感知，尽量以统一方式组织交互与结果展示；其二，后端需要集中处理鉴权封装、参数适配、结果标准化和异步调度；其三，本地模型与外部接口的接入方式差异明显，需要在系统边界上进行隔离。

这种三层结构本质上是一种面向研究原型的松耦合设计。它既能支持本地模型复现，又能保留调用外部能力的灵活性；既能满足实验中的方法对比需求，也能支撑功能扩展和部署演示。因而，该架构与本文“研究验证 + 工程实现”并重的课题定位是匹配的。

根据上述设计思路，本文平台总体架构如图~\ref{fig:system-architecture} 所示。该图并不强调某一个单独模型的内部结构，而是突出前端页面、后端代理、本地模型服务、外部能力接口与样本资源之间的协作关系。这样处理更符合本文系统设计型工作的重点，即说明多源能力如何被统一组织到同一安全分析平台中。

\begin{figure}[H]
\centering
\resizebox{\textwidth}{!}{%
\begin{tikzpicture}[
  font=\small,
  node distance=8mm and 8mm,
  archbox/.style={draw, rounded corners=3pt, thick, align=center, minimum width=2.9cm, minimum height=0.85cm, fill=white},
  groupbox/.style={draw, rounded corners=5pt, thick, inner sep=7pt},
  arrow/.style={-Latex, thick},
  dashedarrow/.style={-Latex, thick, dashed}
]
  \node[archbox, fill=blue!7] (home) at (0,0) {首页与导航};
  \node[archbox, fill=blue!7, right=of home] (genpage) {生成页面};
  \node[archbox, fill=blue!7, right=of genpage] (fakepage) {虚假检测页面};
  \node[archbox, fill=blue!7, right=of fakepage] (safepage) {安全审核页面};
  \node[archbox, fill=blue!7, right=of safepage] (datapage) {内容管理页面};
  \node[groupbox, fit=(home)(genpage)(fakepage)(safepage)(datapage), label={[font=\bfseries]above:前端展示层}] (frontend) {};

  \node[archbox, fill=green!8, below=1.35cm of genpage] (router) {统一API路由};
  \node[archbox, fill=green!8, right=of router] (format) {参数校验\\格式转换};
  \node[archbox, fill=green!8, right=of format] (auth) {密钥隔离\\签名鉴权};
  \node[archbox, fill=green!8, right=of auth] (normalize) {任务调度\\结果归一化};
  \node[groupbox, fit=(router)(format)(auth)(normalize), label={[font=\bfseries]above:后端代理层}] (backend) {};

  \node[archbox, fill=orange!10, below=1.35cm of router] (ufd) {UniversalFakeDetect\\图像伪造检测};
  \node[archbox, fill=orange!10, right=of ufd] (sd) {Stable Diffusion\\本地文生图};
  \node[archbox, fill=orange!10, right=of sd] (ms) {ModelScope T2V\\本地文生视频};
  \node[archbox, fill=orange!10, right=of ms] (volc) {火山方舟/视觉智能\\生成与审核接口};
  \node[groupbox, fit=(ufd)(sd)(ms)(volc), label={[font=\bfseries]above:模型推理与外部能力层}] (capability) {};

  \node[archbox, fill=gray!12, below=1.15cm of sd] (dataset) {公开样本\\\texttt{datasets/}};
  \node[archbox, fill=gray!12, right=of dataset] (mockres) {演示资源\\\texttt{public/mock/}};
  \node[archbox, fill=gray!12, right=of mockres] (records) {生成样本库\\检测记录库};
  \node[groupbox, fit=(dataset)(mockres)(records), label={[font=\bfseries]above:样本与结果资源}] (datares) {};

  \draw[arrow] (genpage.south) -- (router.north);
  \draw[arrow] (fakepage.south) -- (format.north);
  \draw[arrow] (safepage.south) -- (auth.north);
  \draw[arrow] (datapage.south) -- (normalize.north);
  \draw[arrow] (backend.south) -- (capability.north);
  \draw[dashedarrow] (capability.south) -- (datares.north);
  \draw[dashedarrow] (datares.west) .. controls +(0,1.6) and +(0,-1.6) .. (frontend.west);
\end{tikzpicture}}
\caption{AIGC安全分析平台总体架构}
\label{fig:system-architecture}
\end{figure}

\subsection{前端展示层设计}

前端展示层是用户与平台交互的入口，主要承担内容上传、任务选择、检测触发、结果展示、样本保存、一键送检和历史查看等功能。从整体结构看，前端不负责算法推理与复杂业务编排，而是通过相对统一的组件结构组织不同页面。按照当前课题目标，前端可以划分为首页、内容生成页、虚假检测页、不安全检测页和内容管理页等若干页面模块，各页面围绕具体业务场景提供相应操作入口。

为了提高系统一致性，前端展示层的设计原则是“统一入口、分任务展示、结果结构化”。即不同业务任务虽然在具体内容上存在差异，但用户都应遵循相近的操作流程：输入内容、选择任务、触发检测或生成、查看结果。这种统一交互模式能够降低系统复杂度，并增强平台的可演示性。

\subsection{后端代理层设计}

后端代理层是整个系统的核心枢纽。它位于前端与模型服务/外部接口之间，负责统一接收前端请求，并将其转发至正确的能力来源。更重要的是，该层承担了多源能力接入过程中的差异消解工作。例如，不同接口可能在鉴权方式、输入参数形式、异步处理机制和结果返回结构上存在明显不同，如果由前端直接处理这些差异，将导致页面逻辑复杂且难以维护。通过后端代理层统一封装后，前端只需面向统一业务接口即可。

在总体设计中，后端代理层需要支持三类核心任务：第一，统一请求入口，包括参数校验、内容格式转换和任务路由；第二，统一安全控制，包括密钥管理、签名计算和错误屏蔽；第三，统一结果结构，包括真假判断、风险标签、生成结果和历史记录的规范化表示。

\subsection{模型推理与外部能力层设计}

模型推理与外部能力层是系统真正执行检测与生成任务的能力来源，可进一步分为本地模型服务和外部接口服务两类。本地模型服务主要承担可控、可复现的研究任务，如本地图像虚假检测、本地文生图和本地文生视频等；外部接口服务则主要承担开放场景下的图像AI鉴别、视频AI鉴别以及图像/视频安全审核等任务。

这种“双来源能力层”设计与单一后端模型方案相比更符合本文课题实际。一方面，本地模型能够体现毕业设计在算法复现和自主部署方面的工作量；另一方面，外部接口能够弥补本地模型在多模态能力、视频审核能力和工程成熟度方面的不足。对平台而言，关键并不在于能力来源是否统一，而在于能否通过后端代理将这些异构能力组织为一个统一系统。

\section{系统核心业务流程}

\subsection{内容检测流程设计}

系统的基本检测流程可概括为“输入内容—任务选择—请求转发—结果解析—风险展示”。首先，用户在前端选择输入内容并指定任务类型，例如图像虚假检测、视频虚假检测、图像安全审核或视频安全审核。随后，前端将统一格式的请求发送给后端代理层。后端根据任务类型和能力配置，将请求转发至对应的本地模型服务或外部接口，并在收到结果后完成字段标准化与状态归一化处理。最后，前端根据统一结果结构完成风险标签、置信度、原因说明和预览信息的展示。

为了使上述过程更加清晰，本文将内容检测流程抽象为图~\ref{fig:detection-flow}。该流程图体现了平台设计中的两个关键点：一是前端只面向统一任务入口，二是本地模型与外部接口均由后端代理层统一封装，最终以标准化结果返回前端展示。

\begin{figure}[H]
\centering
\resizebox{0.98\textwidth}{!}{%
\begin{tikzpicture}[
  font=\small,
  node distance=8mm and 10mm,
  flowstep/.style={draw, rounded corners=3pt, thick, align=center, minimum width=2.75cm, minimum height=0.9cm, fill=blue!6},
  decision/.style={draw, diamond, aspect=2.0, thick, align=center, inner sep=1pt, fill=yellow!15},
  source/.style={draw, rounded corners=3pt, thick, align=center, minimum width=2.85cm, minimum height=0.9cm, fill=orange!10},
  arrow/.style={-Latex, thick}
]
  \node[flowstep] (input) {图像/视频输入\\文件或URL};
  \node[flowstep, right=of input] (frontend) {前端预处理\\预览与任务选择};
  \node[flowstep, right=of frontend] (backend) {后端代理\\参数校验与路由};
  \node[decision, right=of backend] (dispatch) {能力\\来源};
  \node[source, above right=7mm and 10mm of dispatch] (local) {本地模型服务\\UFD/SD/T2V};
  \node[source, below right=7mm and 10mm of dispatch] (external) {外部平台接口\\方舟/视觉智能};
  \node[flowstep, right=4.2cm of dispatch] (normalize) {结果解析\\字段标准化};
  \node[flowstep, right=of normalize] (visual) {风险可视化\\标签/置信度/说明};
  \node[flowstep, below=1.05cm of visual] (record) {结果管理\\记录与回溯};

  \draw[arrow] (input) -- (frontend);
  \draw[arrow] (frontend) -- (backend);
  \draw[arrow] (backend) -- (dispatch);
  \draw[arrow] (dispatch) -- (local);
  \draw[arrow] (dispatch) -- (external);
  \draw[arrow] (local.east) -| (normalize.north);
  \draw[arrow] (external.east) -| (normalize.south);
  \draw[arrow] (normalize) -- (visual);
  \draw[arrow] (visual) -- (record);
  \draw[arrow] (record.west) .. controls +(-2.0,0) and +(0,-1.1) .. (frontend.south);
\end{tikzpicture}}
\caption{内容安全检测业务流程}
\label{fig:detection-flow}
\end{figure}

这一流程的关键在于“结果解析”环节。由于不同检测后端返回的数据结构差异较大，如果没有统一结果结构，前端就必须为每种能力编写单独渲染逻辑，系统复杂度会迅速增加。因此，在总体设计中，结果结构的统一性是业务流程顺畅运行的核心前提。

\subsection{生成与闭环验证流程设计}

除检测流程外，平台还需要支持“生成—入库—送检—记录”的闭环验证流程。其基本逻辑是：用户先通过生成模块产生图像或视频内容，系统将结果保存为带有来源信息的样本；用户随后可选择真实性检测或安全审核，平台根据样本类型自动进入对应检测页，并带入媒体地址、文件名、模型名称和样本编号；检测完成后，系统再把检测结论写入记录库并更新样本检测状态。该流程在论文中的意义主要体现在两个方面。第一，它为实验样本构建提供平台内闭环来源，减少对外部数据的完全依赖。第二，它能够直观展示系统具备从生成到分析再到回溯的完整链路，有利于体现平台整体性。

在总体设计层面，闭环验证流程并不要求所有生成结果都自动进入检测模块，而是要求平台在数据流上支持这种联动关系。也就是说，生成模块与检测模块之间应当通过统一样本结构共享输出数据，避免重复下载、重复上传和手工登记。图~\ref{fig:generate-detect-loop} 给出了生成结果进入检测链路的设计流程。

\begin{figure}[H]
\centering
\resizebox{0.96\textwidth}{!}{%
\begin{tikzpicture}[
  font=\small,
  node distance=8mm and 10mm,
  loopstep/.style={draw, rounded corners=3pt, thick, align=center, minimum width=2.75cm, minimum height=0.9cm, fill=blue!6},
  storebox/.style={draw, rounded corners=3pt, thick, align=center, minimum width=2.95cm, minimum height=0.9cm, fill=green!8},
  arrow/.style={-Latex, thick}
]
  \node[loopstep] (generate) {生成模块\\产出图像/视频};
  \node[storebox, right=of generate] (sample) {生成样本入库\\保存媒体与元数据};
  \node[loopstep, right=of sample] (choose) {一键送检\\选择检测目标};
  \node[loopstep, right=of choose] (detect) {检测/审核页面\\自动填充输入};
  \node[storebox, right=of detect] (record) {检测记录入库\\更新样本状态};
  \node[loopstep, below=1.05cm of record] (review) {内容管理中心\\回看样本与历史};

  \draw[arrow] (generate) -- (sample);
  \draw[arrow] (sample) -- (choose);
  \draw[arrow] (choose) -- (detect);
  \draw[arrow] (detect) -- (record);
  \draw[arrow] (record) -- (review);
  \draw[arrow] (sample.south) .. controls +(0,-1.0) and +(-3.0,0) .. (review.west);
\end{tikzpicture}}
\caption{生成样本入库与一键送检闭环流程}
\label{fig:generate-detect-loop}
\end{figure}

\subsection{结果管理与追踪流程设计}

为了支撑论文实验与系统演示，平台需要具备结果管理与追踪流程。对于检测任务，应记录任务类型、输入来源、检测结论、风险标签、时间戳和必要的说明字段；对于生成任务，应记录生成类型、生成结果地址、生成时间、来源模块和相关参数摘要。这样做的目的是将一次性的页面交互结果沉淀为可回溯的数据记录。

根据当前平台定位，结果管理并不追求完整生产级资产系统，而是强调研究样本的轻量持久化与关联追踪。表~\ref{tab:content-management-objects} 给出了内容管理模块需要维护的两类核心对象。生成样本侧重“从哪里来、是什么内容、是否已经检测”，检测记录侧重“检测了什么、结论是什么、由哪个样本触发”。二者通过样本编号建立关联，使用户既可以从样本查看检测历史，也可以从检测记录回溯生成来源。

\begin{table}[H]
\centering
\caption{内容管理核心对象设计}
\label{tab:content-management-objects}
\setlength{\tabcolsep}{4pt}
\renewcommand{\arraystretch}{1.35}
\zihao{5}\songti
\begin{tabular}{C{2.5cm}L{5.2cm}L{6.2cm}}
\toprule
\textbf{对象类型} &
\multicolumn{1}{C{5.2cm}}{\textbf{关键字段}} &
\multicolumn{1}{C{6.2cm}}{\textbf{追踪意义}} \\
\midrule
生成样本 & 样本编号、媒体类型、标题、来源模块、生成模型、提示词、媒体地址、存储类型、检测状态、最近检测结果 & 保存平台生成的图像或视频，使其可以被再次预览、下载、送检和复核。 \\
检测记录 & 记录编号、关联样本编号、检测类型、媒体类型、结论、置信度或风险分数、检测模型、标签、原始结果、创建时间 & 保存每次检测或审核结论，使实验结果能够跨页面保留，并支持从记录反查样本来源。 \\
\bottomrule
\end{tabular}
\end{table}

结果追踪流程的引入有助于解决两个问题：一是便于实验统计与人工核查，二是便于回看不同任务的示例结果。即使后续将轻量文件存储替换为数据库，样本对象、检测记录对象和二者关联关系仍可保持稳定。

\section{模块划分与接口关系设计}

\subsection{功能模块划分}

结合上述需求与业务流程，本文将平台划分为内容生成模块、虚假内容检测模块、不安全内容审核模块、内容管理模块和系统支撑模块五类。内容生成模块负责产出研究样本和演示素材，并提供保存与送检入口；虚假内容检测模块负责判断内容真伪与AI生成属性；不安全内容审核模块负责识别内容风险类别与处置建议；内容管理模块负责样本库、检测记录、关联状态和结果回看；系统支撑模块则主要承担鉴权、代理、调度、日志和部署等功能。

这种模块划分方式的优点在于能够对应真实业务任务，便于用户理解系统结构，也便于在工程实现中分别维护生成、检测、审核和结果管理等功能。不同模块之间通过统一输入输出结构进行连接，既避免了功能耦合过重，也为后续接入新模型或扩展新任务保留了空间。

\subsection{统一输入输出设计}

在多源能力接入场景下，统一输入输出设计是降低系统复杂度的关键。对于输入层，平台应尽量将图像与视频内容统一抽象为“文件输入”与“链接输入”两种基本形式，再由后端根据任务需要完成Base64编码、URL透传或其他格式转换。对于输出层，平台则应将不同后端返回的复杂结果归一化为若干核心字段，例如任务类型、判定结果、置信度、风险标签、说明文本、媒体预览地址等。

统一输入输出设计的重要价值在于，它使系统能够在不改变前端主流程的前提下替换底层能力。例如，当后续新增一个图像虚假检测模型时，只要该模型能够在后端层被适配为统一输出结构，前端页面就无需整体重写。这种设计思路也是平台可扩展性的直接体现。

根据当前平台的任务划分，本文将检测类结果抽象为表~\ref{tab:unified-result-schema} 所示的核心字段。该表并不对应某一个具体接口的原始返回，而是后端代理层归一化后的平台内部结果结构。通过这种方式，前端能够用相同的展示组件处理本地模型、多模态理解接口和内容审核接口返回的结果。

\begin{table}[H]
\centering
\caption{检测结果统一字段设计}
\label{tab:unified-result-schema}
\renewcommand{\arraystretch}{1.35}
\zihao{5}\songti
\begin{tabular}{C{2.0cm}C{4.1cm}L{8.0cm}}
\toprule
\textbf{字段类别} &
\textbf{代表字段} &
\multicolumn{1}{C{8.0cm}}{\textbf{设计说明}} \\
\midrule
任务标识 & taskType、mediaType & 标记任务属于虚假检测或安全审核，以及输入媒介为图像或视频。 \\
判定结论 & isFake、safe、suggestion & 虚假检测关注是否为AI生成或伪造；安全审核关注是否安全以及通过、复审、阻断等建议。 \\
可信程度 & confidence、riskScore & 将模型分数或审核建议映射为0到1的置信度或0到100的风险分值，便于统一可视化。 \\
风险解释 & labels、reason、artifacts & 保存风险标签、模型说明和可疑痕迹描述，提升检测结果的可解释性。 \\
展示信息 & previewUrl、model、details & 记录预览资源、能力来源和任务细节，支持页面展示与实验回溯。 \\
追踪信息 & sampleId、sourceModule、createdAt & 关联生成样本、来源模块和检测时间，使检测记录能够回溯到具体生成任务。 \\
\bottomrule
\end{tabular}
\end{table}

在风险等级映射上，本文采用相对克制的三档表达：低风险表示未发现明显问题或建议通过；中风险表示模型认为存在不确定因素，需要人工复审；高风险表示检测结果中存在较强伪造或违规信号，建议阻断或重点核查。这样设计的目的不是替代人工审核结论，而是把不同模型输出转换为一致的工程表达，使用户能够快速理解检测结果的严重程度。

\subsection{模块间协同关系}

各功能模块并不是孤立运行的页面集合，而是围绕统一业务流程协同工作。内容生成模块负责产生图像或视频样本，并将结果保存为可追踪样本；虚假内容检测模块和不安全内容审核模块分别从真实性和安全性两个角度进行分析；内容管理模块则负责沉淀生成样本、检测记录和关联状态。系统支撑模块贯穿整个流程，承担接口代理、鉴权封装、任务调度、错误处理、结果标准化和轻量持久化等基础能力。

在该协同关系中，后端代理层是连接各模块的关键。前端页面只需按照任务类型提交统一请求，后端根据任务配置选择本地模型或外部接口，并将返回结果映射为统一字段。通过这种方式，生成、检测、审核和记录管理能够共享同一套输入输出规范，平台也具备继续扩展新任务的能力。

\section{本章小结}

本章围绕多模态AIGC安全分析平台的建设目标，对系统需求和总体设计进行了系统梳理。首先，从课题定位出发，明确了平台需要兼顾研究验证、工程演示和后续扩展三重目标；其次，围绕图像与视频输入、虚假内容检测、不安全内容审核、内容生成与结果管理等环节，对系统功能需求进行了分析；随后，从可扩展性、可维护性、响应性能以及安全性与可部署性等方面，进一步归纳了平台必须满足的非功能需求。

在此基础上，本文提出了由前端展示层、后端代理层和模型推理与外部能力层构成的三层总体架构，并分析了内容检测流程、生成样本入库与一键送检流程、结果追踪流程，完成了平台功能模块与统一输入输出结构的总体设计。通过这一章的分析可以看出，本文系统并非简单堆叠多个模型接口，而是在明确需求约束的前提下，构建了一个具有统一组织逻辑、样本追踪能力和扩展能力的研究原型平台。
